\section{Restriction and cohomology bounds}
\label{sec:analysis}

The estimates in this section convert numerical control of slope factors
into control of sections and first cohomology.  The constants depend on the
polarized variety and on specified twist ranges, but not on the ranks of
the sheaves.  This distinction is essential: the first-cohomology estimate
will be proportional to a discriminant error, which can vanish even when
the rank is large.

Throughout this section, $H$ is very ample and $h=H^3$.  On a smooth
surface $S\in|H|$, we use
\[
 d_T(G)=\frac{H|_S\,\ch_1^T(G)}h,
 \qquad w_T(G)=\frac{\ch_2^T(G)}h,
 \qquad \mu_T(G)=\frac{d_T(G)}{\rk G}.
\]
For divisor classes on $S$, put $\langle C,D\rangle=CD/h$ and use the
norm
\[
 \|C\|^2=\langle H,C\rangle^2-\langle C_\perp,C_\perp\rangle,
 \qquad C_\perp=C-\langle H,C\rangle H.
\]
Here and below $H$ also denotes its restriction.  The definitions of
$\xi_G$ and $\delta_G^T$ consequently have the same formulas as on $X$.
In particular, for a positive-rank sheaf they give the identity
\begin{equation}\label{eq:surface-norm-discriminant}
 \mu_T(G)^2+\delta_G^T
 =\left\|\frac{c_1(G)}{\rk G}-T\right\|^2+\xi_G.
\end{equation}
On a smooth curve $C_m\in|mH|_S$, normalize the slope as
$\mu(G)=\deg(G)/(mh\rk G)$.  This makes the slope of an unmodified
restriction equal to the slope before restriction.

The first-cohomology estimate will concern stable surface sheaves
whose normalized first character is close to a negative scalar class.
For the fixed rational twist $B_0$, write
\[
 T=B_0|_S+(b-A)H,\qquad |b|\le1.
\]
For sufficiently large $A$, our target is to bound first cohomology
at fixed twists by sums of
\[
 r(G)\left(\left\|\frac{c_1(G)}{r(G)}-T\right\|^2+\xi_G\right).
\]
This quantity can vanish even when the total rank is large.  The
negative scalar part of $T$ will give positive central curvature on
the dual bundle; its trace-free curvature will then control the
remaining harmonic forms.  A separate estimate for sections uses the
rank of a subsheaf's evaluation map.  We first choose surfaces and
curves with uniform geometry, so that both estimates have constants
independent of the sheaves and their ranks.

\subsection{Restriction and the choice of smooth sections}

\begin{lemma}[Restriction variance]\label{lem:restriction-variance}
Let $G$ be a slope-semistable torsion-free sheaf of rank $r>0$ on $X$.
For a very general smooth $S\in|H|$, restrict $G$ and refine its slope
Harder--Narasimhan filtration into slope-stable torsion-free factors
$G_i$, with ranks $r_i$ and normalized slopes $\mu_i$.  Then
\begin{equation}\label{eq:restriction-variance}
 \sum_i r_i\bigl(\mu_i-\mu(G)\bigr)^2\le r\xi_G.
\end{equation}
The same assertion holds for a slope-semistable torsion-free sheaf on
$S$, restricted to a very general smooth $C_m\in|mH|_S$, for every
positive integer $m$.

More generally, restrict a finite filtration with semistable
torsion-free factors and refine each restricted factor separately.
From $X$ to $S$, for any fixed divisor twist $T$ on $X$, restricted
to $S$, the increases in
$\sum_i r_i\mu_T(G_i)^2$ and in $\sum_i r_i\delta_{G_i}^T$ are equal
and lie between zero and the original sum $\sum_G r(G)\xi_G$.
From $S$ to $C_m$, the same upper bound holds for the increase of the
second slope moment, using any fixed divisor twist on $S$ and its
restriction to $C_m$.  The concatenated filtration need not be the
Harder--Narasimhan filtration of the whole restriction.
\end{lemma}

\begin{proof}
Write
\[
 \operatorname{Disc}(G)
 =2r c_2(G)-(r-1)c_1(G)^2=c_1(G)^2-2r\ch_2(G).
\]
The characteristic-zero semistable case of Langer's restriction
inequality \cite[Theorem~3.1]{Langer2004} has zero extremal-slope
term.  On $X$ it gives, for raw restricted slopes $s_i=h\mu_i$,
\[
 \sum_{i<j}r_i r_j(s_i-s_j)^2
 \le hH\operatorname{Disc}(G)=h^2r^2\xi_G.
\]
Since $\sum r_i\mu_i=r\mu(G)$, the left side is
$h^2r\sum_i r_i(\mu_i-\mu(G))^2$.
Division gives \eqref{eq:restriction-variance}.  If the restriction is
semistable, its left side is zero, and the inequality follows directly
from $\xi_G\ge0$.  Refinement at a fixed slope leaves the sum unchanged.
For a surface restriction, apply the same theorem with its very ample
divisor $mH$.  The raw slopes are $s_i=mh\mu_i$ and the right side is
$m^2h\operatorname{Disc}(G)=m^2h^2r^2\xi_G$; the factors $m^2h^2r$
cancel in the same way.

We use the corrected very-general formulation of the theorem
\cite[p.~1211]{Langer2004Addendum}.  In the ample polarizations used
here, that addendum also permits general divisors: the stated
difficulty concerns the remaining nef polarizations when they are
not ample.

For each original factor, the weighted mean of its restricted slopes
is unchanged.  Subtracting a fixed scalar $\mu(T)$ from these slopes
therefore does not change their variance.  The elementary identity
\[
 \sum_i r_i(\mu_i-\mu(T))^2
 =r(\mu(G)-\mu(T))^2
   +\sum_i r_i(\mu_i-\mu(G))^2
\]
proves the second-moment assertion.  From $X$ to $S$, restriction
preserves the total $w_T$; hence
$\sum r_i\delta_{G_i}^T=\sum r_i\mu_T(G_i)^2-2w_T(G)$ changes by
exactly the same amount.  Sum these identities over the original
factors.  Their restricted and refined filtrations are exact for a
suitable choice of the divisor, as justified next.
\end{proof}

\begin{lemma}[Uniform smooth sections]\label{lem:uniform-sections}
There is an integer $m>0$, depending only on $X,H$, such that every
smooth $S\in|H|$ satisfies
\begin{equation}\label{eq:surface-fixed-m}
 h^0(S,\OO_S(mH))=\chi(\OO_S(mH)),
 \qquad \mu(K_S)-m\le-1,
\end{equation}
and $H^0(X,\OO_X(mH))\to H^0(S,\OO_S(mH))$ is surjective.
One can choose relatively compact analytic neighborhoods of a smooth
transverse pair of divisors in $|H|\times|mH|$ such that the resulting
surfaces and curves, equipped with the induced Fubini--Study metrics,
have uniform smooth local geometry.

Within these neighborhoods one may impose, for each application, the
very-general restriction conditions of
Lemma~\ref{lem:restriction-variance} and exactness of any specified
finite collection of restricted sequences.  Restrictions of a
specified finite collection of torsion-free sheaves may also be
required to remain torsion-free.  For a specified bounded complex,
one may require its derived restriction to have ordinary cohomology
equal to the restrictions of its ordinary cohomology sheaves.
\end{lemma}

\begin{proof}
Choose $m$ so large that $(m-1)H-K_X$ is ample,
$H^1(X,\OO_X((m-1)H))=0$, and $m\ge\mu(K_S)+1$.
The last slope is independent of $S$, by adjunction
$K_S=(K_X+H)|_S$.  Kodaira vanishing on $S$
\cite[Chapter~VII, Theorem~3.3]{DemaillyCADG} gives the first equality
in \eqref{eq:surface-fixed-m}, and the restriction sequence on $X$
gives surjectivity on sections.

Bertini's theorem supplies a smooth pair meeting transversely.
Shrink an analytic neighborhood of that pair so that its closure is
compact and all intersections remain smooth.  The universal
surfaces and curves over this closure are smooth proper families.
Local trivializations on a finite cover give uniform bounds for
the metrics, their derivatives, coordinate comparisons, and the
partitions of unity used below.  Surjectivity on sections allows the
curve on each chosen surface to vary in an analytic open subset of
its complete linear system.

For a fixed coherent sheaf, a general defining section avoids all
its associated points, and is thus a nonzerodivisor.  Applying this
to the quotients in a finite collection of exact sequences preserves
their exactness under restriction.  To ensure torsion-freeness,
embed a torsion-free sheaf into a vector bundle and require the
defining section to be a nonzerodivisor on its cokernel.  Restriction
then embeds the restricted sheaf into the restricted bundle.  Such
an embedding exists because a sufficiently positive twist of the
dual is globally generated and the original sheaf injects into its
double dual.  Repeating the argument on the surface gives the curve
assertion.  Applied to the ordinary cohomology sheaves of a bounded
complex, the nonzerodivisor condition kills the first Tor terms in
the restriction spectral sequence, giving the asserted cohomology
identification.

These are finitely many nonempty Zariski-open conditions at each
stage.  A very-general condition removes at most countably many
proper algebraic closed subsets.  Each such subset has empty
analytic interior, so the Baire theorem ensures a choice in every
one of our nonempty analytic neighborhoods.  The geometric
constants were fixed before this choice and do not depend on the
sheaves being restricted.
\end{proof}

\subsection{Sections of subsheaves}

We first bound sections using a scalar heat estimate.  The rank of a
subsheaf enters only through the rank of its evaluation map; the
ambient bundle may have arbitrarily large rank.

\begin{lemma}[Uniform heat estimates]\label{lem:uniform-heat}
Let $Y$ range over the compact families of curves or surfaces in
Lemma~\ref{lem:uniform-sections}, with real dimension $n$.  Write
$\Delta_{\mathrm{sc}}$ for the nonnegative scalar Laplacian.  Its heat
kernel satisfies
\[
 0\le k_s(x,y)\le C s^{-n/2}\qquad(0<s\le1).
\]
For any Hermitian vector bundle with metric connection $\nabla$,
put $L=\nabla^*\nabla+1$.  The kernel $K_s(x,y)$ of $e^{-sL}$
satisfies
\begin{equation}\label{eq:connection-heat}
 \|K_s(x,y)\|_{\mathrm{op}}\le C s^{-n/2}
 \qquad(s>0),
\end{equation}
with a constant independent of the bundle, its rank, and its
connection.
\end{lemma}

\begin{proof}
The uniform charts and partitions of unity give the Nash inequality
\begin{equation}\label{eq:nash-uniform}
 \|f\|_2^{2+4/n}
 \le C\bigl(\|df\|_2^2+\|f\|_2^2\bigr)\|f\|_1^{4/n}.
\end{equation}
For completeness, in a Euclidean chart split the Fourier integral
at frequency $R$: its low-frequency part is bounded by
$CR^n\|f\|_1^2$, and its high-frequency part by
$R^{-2}\|df\|_2^2$.  Optimize in $R$, and use a fixed finite
partition of unity.  Derivatives of the partition produce the
additional $\|f\|_2^2$ term.  All comparison constants are uniform
over the chosen families.

Set $L_{\mathrm{sc}}=\Delta_{\mathrm{sc}}+1$ and
$f_s=e^{-sL_{\mathrm{sc}}}f$.
The scalar semigroup is positivity preserving and contracts $L^1$.
For $y(s)=\|f_s\|_2^2$, \eqref{eq:nash-uniform} and
$y'=-2\langle L_{\mathrm{sc}}f_s,f_s\rangle$ imply
\[
 y'\le-c\|f\|_1^{-4/n}y^{1+2/n}.
\]
Integration gives $\|e^{-sL_{\mathrm{sc}}}\|_{1\to2}\le Cs^{-n/4}$.
Selfadjointness and composition give
$\|e^{-sL_{\mathrm{sc}}}\|_{1\to\infty}\le Cs^{-n/2}$ for every $s>0$.
Multiplication by $e^s$ gives the stated estimate for $k_s$ when
$s\le1$.

For a section-valued heat solution, Kato's inequality and the
maximum principle give
\[
 |e^{-sL}v|(x)\le(e^{-sL_{\mathrm{sc}}}|v|)(x).
\]
One obtains the differential inequality by differentiating
$(|v|^2+\varepsilon^2)^{1/2}$ and then letting
$\varepsilon\downarrow0$; compatibility of the connection with the
metric is the only bundle input.  Approximate a vector-valued point
mass at $y$ in this inequality.  This bounds the operator norm of
$K_s(x,y)$ by the scalar kernel of $e^{-sL_{\mathrm{sc}}}$, and proves
\eqref{eq:connection-heat}.
\end{proof}

\begin{lemma}[Sections of a subsheaf]\label{lem:subsheaf-sections}
Let $Y$ be one of the curves or surfaces in
Lemma~\ref{lem:uniform-sections}, let $G$ be a slope-stable
torsion-free sheaf on $Y$, and let $J\subset G$ have rank $q$.
Then
\begin{equation}\label{eq:subsheaf-sections}
 h^0(Y,J)\le Cq\bigl(1+\mu(G)_+\bigr)^{\dim Y},
 \qquad x_+=\max\{x,0\},
\end{equation}
where $C$ depends only on the chosen geometric families.  The same
bound applies to subsheaves of a stable dual bundle tensored with
a line bundle, using the slope of that twisted bundle.
\end{lemma}

\begin{proof}
On a curve $G$ is locally free.  On a surface its reflexive hull
$W=G^{**}$ is locally free and slope-stable: intersecting a
subsheaf of $W$ with $G$ preserves rank and first Chern class, since
$W/G$ is zero-dimensional.  The Hermitian--Einstein theorem
\cite{Donaldson1985,UhlenbeckYau1986} therefore gives a metric on
$W$ for which $i\Lambda F_W$ is the scalar determined by
$\mu(G)$.  Its proportionality constant is fixed by the dimension
and our volume normalization.

For a holomorphic section $s$ of $W$, the Bochner identity gives
\[
 \Delta_{\mathrm{sc}}|s|^2\le C\mu(G)_+|s|^2.
\]
Scalar heat comparison consequently yields
$|s(x)|^2\le e^{C u\mu(G)_+}(e^{-u\Delta_{\mathrm{sc}}}|s|^2)(x)$.
Take $u=(1+\mu(G)_+)^{-1}$ and apply
Lemma~\ref{lem:uniform-heat}.  The evaluation operator, with the
$L^2$ norm on $H^0(Y,W)$, satisfies
\begin{equation}\label{eq:evaluation-operator}
 \|\operatorname{ev}_x\|_{\mathrm{op}}^2
 \le C\bigl(1+\mu(G)_+\bigr)^{\dim Y}.
\end{equation}
Give $H^0(Y,J)\subset H^0(Y,W)$ the induced $L^2$ norm and choose
an orthonormal basis $s_1,\ldots,s_N$.  Outside a proper analytic
subset, their evaluations lie in a subspace of dimension at most
$q$.  Thus
\[
 \sum_{j=1}^N|s_j(x)|^2
 =\|\operatorname{ev}_x|_{H^0(J)}\|_{\mathrm{HS}}^2
 \le q\|\operatorname{ev}_x\|_{\mathrm{op}}^2.
\]
Integration proves \eqref{eq:subsheaf-sections}; the volumes are
uniformly bounded.  The case $q=0$ means $J=0$.  Dualization and
tensoring with a line preserve slope stability, so the same proof
gives the last assertion.
\end{proof}

\subsection{A matrix spectral estimate on surfaces}

To control first cohomology, a bound proportional to rank would
not suffice.  We instead count harmonic forms using the square of
the Hilbert--Schmidt norm of the trace-free curvature.  The following
spectral estimate is a four-dimensional, bundle-valued
Cwikel--Lieb--Rozenblum estimate; compare
\cite[Theorem~3.2 and Lemma~3.4]{Frank2014}. We give the
heat-semigroup proof, retaining the matrix trace to keep its constants
independent of rank.

\begin{lemma}[Spectral dimension bound]\label{lem:matrix-spectral}
Let $S$ be one of the surfaces in Lemma~\ref{lem:uniform-sections}.
Let $\mathcal V$ be a Hermitian vector bundle with metric connection,
$L=\nabla^*\nabla+1$, and let $P$ be a bounded measurable
nonnegative selfadjoint endomorphism of $\mathcal V$.
If a finite-dimensional subspace $U$ of the form domain of $L$
satisfies
\begin{equation}\label{eq:spectral-form-hypothesis}
 \langle Lu,u\rangle\le\langle Pu,u\rangle\qquad(u\in U),
\end{equation}
then
\begin{equation}\label{eq:matrix-spectral-bound}
 \dim U\le C\int_S\operatorname{tr}(P^2).
\end{equation}
The constant is independent of $\mathcal V$, its rank, $\nabla$,
and $P$.
\end{lemma}

\begin{proof}
The compact positive operator $L^{-1/2}PL^{-1/2}$ has at least
$\dim U$ eigenvalues greater than or equal to $1$: apply min--max
to the subspace $L^{1/2}U$ and
\eqref{eq:spectral-form-hypothesis}.  Its nonzero eigenvalues,
with multiplicity, are those of $P^{1/2}L^{-1}P^{1/2}$.
Define an operator with an additional time variable by
\[
 (\mathcal Bv)(s)=e^{-sL/2}P^{1/2}v,
 \qquad
 \mathcal B:L^2(S,\mathcal V)\longrightarrow
 L^2((0,\infty)\times S,\mathcal V).
\]
Integration of the heat semigroup gives
$\mathcal B^*\mathcal B=P^{1/2}L^{-1}P^{1/2}$.

Fix $0<\eta<1$ and split $\mathcal B=\mathcal B_{\le}+\mathcal B_>$
by inserting the fiberwise spectral projections
$\mathbf1_{sP\le\eta}$ and $\mathbf1_{sP>\eta}$ immediately after
$P^{1/2}$.  Heat contraction gives
\[
 \|\mathcal B_{\le}v\|^2
 \le\int_0^\infty\langle P\mathbf1_{sP\le\eta}v,v\rangle\,ds
 \le\eta\|v\|^2.
\]
Indeed, each positive eigenvalue $p$ contributes
$\int_0^{\eta/p}p\,ds=\eta$, and a zero eigenvalue contributes zero.

Let $K_s$ be the kernel of $e^{-sL}$.  The semigroup identity and
Lemma~\ref{lem:uniform-heat}, in real dimension four, imply
\begin{align*}
 \|\mathcal B_>\|_{\mathrm{HS}}^2
 &=\int_0^\infty\int_S
   \operatorname{tr}\bigl(P\mathbf1_{sP>\eta}K_s(x,x)\bigr)\,dx\,ds\\
 &\le C\int_S\int_0^\infty
   s^{-2}\operatorname{tr}\bigl(P\mathbf1_{sP>\eta}\bigr)\,ds\,dx\\
 &=C\eta^{-1}\int_S\operatorname{tr}(P^2).
\end{align*}
The last identity follows eigenvalue by eigenvalue from
$\int_{\eta/p}^\infty s^{-2}p\,ds=p^2/\eta$.
All integrands are nonnegative, so Tonelli's theorem applies.

On the spectral subspace where $\mathcal B^*\mathcal B\ge1$,
the triangle inequality gives
$\|\mathcal B_>v\|\ge(1-\sqrt\eta)\|v\|$.
Apply this to an orthonormal basis of that subspace and sum.
Taking $\eta=1/4$ proves \eqref{eq:matrix-spectral-bound}.
\end{proof}

In the application below, $P$ will be built from trace-free curvature.
The matrix trace in \eqref{eq:matrix-spectral-bound} will therefore
bound the dimension of harmonic forms by curvature energy, with no
additional term proportional to the bundle rank.

\subsection{First cohomology near a negative scalar class}

Fix the rational class $B_0$ from the threefold, and restrict it to
the chosen surfaces.  The next proposition concerns sheaves whose
normalized first Chern classes are close to
$B_0+(b-A)H$.  A sufficiently negative scalar part makes the central
curvature component of the dual bundle positive; the trace-free curvature measures the
failure of the resulting vanishing theorem.

\begin{proposition}[Surface first cohomology]\label{prop:negative-surface-cohomology}
Fix a real number $M$.  There is a number $A_0$, depending only on
$X,H,B_0,M$ and the chosen smooth neighborhoods, with the following
property.  For $A\ge A_0$ and a finite set
$I\subset\{t\in\ZZ:t\le M\}$, there is a constant $C$ such that,
for every chosen surface $S$, every $b\in[-1,1]$, every finite
collection of slope-stable torsion-free sheaves $G_i$ on $S$, and
every $t\in I$,
\begin{equation}\label{eq:negative-surface-cohomology}
 \sum_i h^1(S,G_i(tH))
 \le C\sum_i r_i\left(
  \left\|\frac{c_1(G_i)}{r_i}-T\right\|^2+\xi_{G_i}\right),
 \qquad T=B_0|_S+(b-A)H.
\end{equation}
The constant is independent of the collection and its ranks.
The same upper bound holds for $h^1(S,G(tH))$ if the $G_i$ are
the factors of a filtration of $G$.
In particular, a bound $C_0e$ for the sum on the right gives a
bound $CC_0e$ for the first cohomology.
\end{proposition}

The order of constants matters in the next section.  The threshold
$A_0$ uses only the upper twist bound $M$.  Once $A$ and the finite
set $I$ are fixed, the cohomology constant $C$ may depend on both;
it remains independent of the sheaves, their ranks, and $b$.

\begin{proof}
We prove the bound for one factor $G$, writing
$r=\rk G$, $C_1^{\mathrm{num}}=c_1(G)/r$, and
$\rho=\|C_1^{\mathrm{num}}-T\|$.  We separate $\rho\ge1$,
where the numerical error already controls rank, from $\rho<1$,
where the curvature argument is needed.

Suppose first that $\rho\ge1$.  The class $T$ is uniformly bounded
once $A$ is fixed, so $1+\|C_1^{\mathrm{num}}\|^2\le C_A\rho^2$.
Lemma~\ref{lem:subsheaf-sections}, applied to $G(tH)$, gives
$h^0(G(tH))\le Cr(1+\|C_1^{\mathrm{num}}\|^2+t^2)$.
Put $W=G^{**}$.  Serre duality identifies $H^2(G(tH))^*$ with
$\Hom(G,K_S(-tH))=H^0(W^*\otimes K_S(-tH))$; the equality follows
by extending across the finite support of $W/G$.
The dual version of Lemma~\ref{lem:subsheaf-sections} gives the
same bound for $h^2$.

Surface Riemann--Roch reads
\[
 \chi(G(tH))
 =r\chi(\OO_S)+\frac{hr}{2}
   \left(\langle C_1^{\mathrm{num}}+tH,
                    C_1^{\mathrm{num}}+tH\rangle-\xi_G\right)
   -\frac r2(C_1^{\mathrm{num}}+tH)K_S.
\]
The Hodge norm controls both the intersection form and pairing
with $K_S$.  Consequently
$|\chi(G(tH))|\le Cr(1+\|C_1^{\mathrm{num}}\|^2+t^2+\xi_G)$.
Since $I$ is finite, $h^1=h^0+h^2-\chi$ is at most
$Cr(\rho^2+\xi_G)$, as required.

Now assume $\rho<1$.  The hull sequence
$0\to G\to W\to Q\to0$ has $Q$ zero-dimensional and
\begin{equation}\label{eq:hull-length}
 \xi_G-\xi_W=\frac{2\operatorname{length}Q}{hr},
 \qquad
 \operatorname{length}Q\le\frac{hr\xi_G}{2}.
\end{equation}
Here $W$ is stable, as in the proof of
Lemma~\ref{lem:subsheaf-sections}, and $\xi_W\ge0$ by the
classical Bogomolov--Gieseker inequality.  The long exact
cohomology sequence gives
$h^1(G(tH))\le h^1(W(tH))+\operatorname{length}Q$.
It remains to bound $h^1(W(tH))$ by $Cr\xi_W$.

Equip $W$ with its Hermitian--Einstein metric.  Write its Chern
curvature as
\[
 F_W=F_0+\frac{\operatorname{tr}F_W}{r}\operatorname{id}_W,
 \qquad
 \vartheta=\frac{i}{2\pi r}\operatorname{tr}F_W.
\]
The real $(1,1)$-form $\vartheta$ is closed and has constant
contraction, hence is harmonic.  It represents
$C_1^{\mathrm{num}}$.  The Hodge star formula on primitive
$(1,1)$-forms identifies its $L^2$ norm with a fixed multiple of
the divisor norm.  Uniform elliptic estimates on our compact
smooth families therefore give a uniform bound for the supremum
norm of a harmonic representative in terms of this divisor norm.
In particular, since $\rho<1$ and $|b|\le1$,
\begin{equation}\label{eq:central-curvature-positive}
 -\vartheta-t\omega_S\ge(A-M-C_1)\omega_S
 \qquad(t\le M).
\end{equation}
The fixed constant $C_1$ bounds the harmonic representatives of
$B_0|_S+bH$ and the unit ball of errors.  Here $\omega_S$ is the
induced Kähler form, and the line bundle $\OO_S(H)$ has curvature
$2\pi\omega_S$.  Choose $A_0$ so that the right side is strictly
positive with a fixed lower bound.  This choice uses only the
upper endpoint $M$, not the lower endpoint of $I$.

The trace-free curvature $F_0$ is primitive.  Chern--Weil and the
primitive Hodge star formula give
\begin{equation}\label{eq:tracefree-energy}
 \int_S|F_0|_{\mathrm{HS}}^2\le C\frac{\operatorname{Disc}(W)}r
   =Chr\xi_W.
\end{equation}
Indeed
$\operatorname{Disc}(W)=-r\int_S
\operatorname{tr}((iF_0/2\pi)\wedge(iF_0/2\pi))$,
and the last integrand is the negative squared norm because
$F_0$ is primitive.  The fixed conversion constant depends on the
curvature convention, not on $r$.

By Serre duality and the Dolbeault theorem, $h^1(W(tH))$ is the
dimension of the space $U$ of harmonic $(2,1)$-forms with values in
$W^*(-tH)$.  On this bundle the Bochner--Kodaira--Nakano identity
\cite[Chapter~VII, Corollary~1.3]{DemaillyCADG} is
\begin{equation}\label{eq:nakano-identity}
 \Delta''=\Delta'+[iF,\Lambda].
\end{equation}
For the central curvature, \eqref{eq:central-curvature-positive}
makes the commutator on $(2,1)$-forms at least $\lambda I$ for a
fixed $\lambda>0$.  To see the sign directly, diagonalize the
central $(1,1)$-form in a unitary frame.  On top holomorphic
degree, its commutator is the sum of the positive eigenvalues
whose antiholomorphic indices occur; there is one such index
here.  Denote the remaining selfadjoint curvature endomorphism by
$K_0$.  Wedge and contraction act on fixed-dimensional form
spaces, so
\begin{equation}\label{eq:curvature-endomorphisms}
 |K_0|_{\mathrm{HS}}\le C|F_0|_{\mathrm{HS}}.
\end{equation}

We explain explicitly how to retain an elliptic operator in
\eqref{eq:nakano-identity}.  Let $\nabla$ be the total metric
connection on $(2,1)$-forms with these bundle coefficients.
For a pure-type form the terms in $\partial u+\bar\partial u$
have different types, as do their adjoints.  Hence the quadratic
form of the connection de Rham Laplacian is the sum of the forms
of $\Delta'$ and $\Delta''$.  On $U$, the latter vanishes, so the
de Rham form equals the $\Delta'$ form.  The Weitzenböck formula,
restricted to this type, consequently gives
\begin{equation}\label{eq:weitzenbock-on-harmonics}
 \langle\Delta'u,u\rangle
 \ge\|\nabla u\|_2^2-C_2\|u\|_2^2+\langle R_0u,u\rangle
 \qquad(u\in U),
\end{equation}
where $R_0$ is selfadjoint and
$|R_0|_{\mathrm{HS}}\le C|F_0|_{\mathrm{HS}}$.
The bounded term contains the tangent curvature and the central
bundle curvature; its bound $C_2$ is uniform for fixed $A,I$.
The formula follows by anticommuting wedge and contraction in
$d_\nabla$ and $d_\nabla^*$: the second-order part is
$\nabla^*\nabla$, and the remaining terms contract the tangent
and bundle curvature.  Only the fixed form dimension enters
the contraction constants.

Choose
$0<\varepsilon\le\min\{1,\lambda/(C_2+1)\}$.
On $U$, retain $\varepsilon\Delta'$ in
\eqref{eq:nakano-identity}, discard the nonnegative remainder,
and use \eqref{eq:weitzenbock-on-harmonics}.  This gives
\[
 \varepsilon\|\nabla u\|_2^2+
  (\lambda-\varepsilon C_2)\|u\|_2^2
 \le\langle(|K_0|+\varepsilon|R_0|)u,u\rangle.
\]
Our choice ensures $\lambda-\varepsilon C_2\ge\varepsilon$.
Thus, with spectral absolute values of selfadjoint matrices,
\[
 L=\nabla^*\nabla+1,\qquad
 P=\varepsilon^{-1}|K_0|+|R_0|,
 \qquad
 \langle Lu,u\rangle\le\langle Pu,u\rangle\quad(u\in U).
\]
Moreover,
\begin{equation}\label{eq:potential-trace}
 \operatorname{tr}(P^2)
 \le2\varepsilon^{-2}|K_0|_{\mathrm{HS}}^2
       +2|R_0|_{\mathrm{HS}}^2
 \le C|F_0|_{\mathrm{HS}}^2.
\end{equation}
This uses the Hilbert--Schmidt norm of the matrices, without
replacing them by their operator norms times the identity.

Lemma~\ref{lem:matrix-spectral}, followed by
\eqref{eq:tracefree-energy}, now proves
$h^1(W(tH))=\dim U\le Cr\xi_W$.
Together with \eqref{eq:hull-length}, this gives the required bound
for the factors with $\rho<1$ as well.  Sum over the factors.
For an extension $0\to G'\to G\to G''\to0$, the cohomology
sequence gives $h^1(G(tH))\le h^1(G'(tH))+h^1(G''(tH))$;
induction proves the final filtration assertion.
\end{proof}

The combination of Lemma~\ref{lem:restriction-variance} and
Proposition~\ref{prop:negative-surface-cohomology} is the point of
the section.  Restriction increases the relevant numerical sums by
at most a discriminant contribution, while the cohomology bound
uses those sums without adding a bare rank term.  The fixed-apex
argument can therefore apply these estimates even when its
numerical error is zero.
